% Generated draft: requires mathematical and source review.
\documentclass[11pt]{article}
\usepackage[margin=25mm]{geometry}
\usepackage{fontspec}
\setmainfont{Latin Modern Roman}
\usepackage{amsmath,amssymb,mathtools,unicode-math}
\setmathfont{Latin Modern Math}
\usepackage{enumitem,needspace,xurl,hyperref}
\setlength{\parindent}{0pt}
\setlength{\parskip}{5pt}
\setlength{\emergencystretch}{3em}
\newcommand{\N}{\mathbb{N}}
\newcommand{\Z}{\mathbb{Z}}
\newcommand{\Q}{\mathbb{Q}}
\newcommand{\R}{\mathbb{R}}
\newcommand{\C}{\mathbb{C}}
\newcommand{\F}{\mathbb{F}}
\newcommand{\A}{\mathbb{A}}
\newcommand{\PP}{\mathbb{P}}
\newcommand{\E}{\mathbb{E}}
\newcommand{\eps}{\varepsilon}
\newcommand{\dd}{\,\mathrm{d}}
\begin{document}
\section*{UnsolvedMath ID: 30006345}
\textbf{Source problem code:} catalogue entry
\subsection{Definitions and mathematical statement}
Let \(p\) be a prime and let \(k\) be an arbitrary field of characteristic \(p\). A finite \(p\)-group is a finite group whose order is a power of \(p\). A group \(G\) has exponent \(p\) if every element satisfies \(g^p=1\) and some element has order \(p\). Write \(Z(G)\) for its center and \([G,G]\) for the subgroup generated by the commutators \([x,y]=x^{-1}y^{-1}xy\). Nilpotency class exactly two means
\[
\{1\}\ne [G,G]\subseteq Z(G).
\]
The group algebra \(kG\) is the \(k\)-vector space with basis \(G\), with multiplication extending the group multiplication \(k\)-bilinearly.

Determine whether the following assertion holds: for every prime \(p\), every field \(k\) of characteristic \(p\), and every pair of finite \(p\)-groups \(G,H\), both of exponent \(p\) and nilpotency class exactly two,
\[
kG\cong_{k\text{-alg}}kH
\quad\Longrightarrow\quad
G\cong H.
\]
Here the algebra isomorphism is taken to be a unital \(k\)-algebra isomorphism, hence fixes the scalar field. The supplied statement does not explicitly specify the type of algebra isomorphism; this interpretation follows the conventional formulation of the modular isomorphism problem and requires source confirmation.

The supplied report excerpt states that this group class has a positive solution over the prime field \(\F_p\), dating to the 1970s and attributed there to reference [7], while the question over general fields remains open. The bibliographic identity of [7] is not supplied. These historical and status assertions are supplied source qualifications, not independently verified conclusions of this draft.
\subsection{Short English statement}
For finite groups of exponent \(p\) and nilpotency class exactly two, does the group algebra over every field of characteristic \(p\) determine the group up to isomorphism? The supplied source reports a positive answer over \(\F_p\) and describes the arbitrary-field question as open. Algebra isomorphisms are interpreted as fixing the field, subject to source confirmation.
\subsection{Sources}
\begin{itemize}
\item Local UnsolvedMath registry, permanent ID \(30006345\), title: Modular Isomorphism for Class-Two Exponent-\(p\) Groups. Exact source problem code: \texttt{catalogue entry}. Collection: Oberwolfach Reports Open Problems; report year: 2025. The supplied record labels the problem \texttt{open} and the statement \texttt{corrected\_verified}; its verification note reports that the official report was checked and the question isolated and its notation repaired. These are supplied registry claims, not verification performed for this draft.
\item Computational Group Theory (2025), Oberwolfach Reports, pp. 1389--1450, DOI 10.4171/owr/2025/27: \url{https://doi.org/10.4171/owr/2025/27}. The supplied excerpt attributes the presentation ``On the Modular Isomorphism Problem'' to Leo Margolis, with joint work with Taro Sakurai. No separate proposer is identified. Its claim of a positive prime-field solution in the 1970s cites reference [7], whose bibliographic details are absent from the supplied record.
\item Margolis--Sakurai, Modular Isomorphism Problem over all fields: \url{https://arxiv.org/abs/2505.05902}.
\item Modular Isomorphism Problem survey: \url{https://doi.org/10.1365/s13291-022-00249-5}.
\item Oberwolfach Report 27/2025: \url{https://ems.press/content/serial-article-files/51860}.
\item Literature assessment dated exactly 2026-08-22, supplied as dated curation and not independently verified: the modular isomorphism problem has counterexamples in general and many solved subclasses; the exponent-\(p\), class-two question over arbitrary characteristic-\(p\) fields remains unresolved, and recent work extends selected positive cases only. The supplied background describes this assessment as literature triage derived from the OWR curation and requests independent verification of time-sensitive status and cited claims.
\end{itemize}
\end{document}
